\chapter{Teoria przecięć, grupy Chowa i klasy Cherna}
\label{ch:ga7-chow-chern}

\section{Motywacja: dlaczego krzywe musimy liczyć z krotnościami?}
\label{sec:ga7-motivation}

W rozdziale~\ref{ch:ga6-powierzchnie} spotkaliśmy dwie pozornie różne
operacje: przecinanie krzywych i obliczanie stopnia wiązki liniowej.
Są one przejawami tej samej teorii. Jeżeli $C,D$ są krzywymi na
gładkiej powierzchni $S$, a $C\cap D$ jest skończony, chcemy
uzyskać liczbę $C\cdot D$, która uwzględnia styczność,
pozostaje niezmiennicza przy zastępowaniu dzielnika innym
liniowo równoważnym i daje się obliczać bez rysowania
wszystkich punktów przecięcia.

\begin{przyklad}[Jedno przecięcie, lecz liczba większa od jedności]
Dla krzywych $C=\V(y)$ i $D=\V(y-x^3)$ w $\A^2_k$ jest
$C\cap D=\{(0,0)\}$, ale
$k[x,y]_{(x,y)}/(y,y-x^3)\cong k[x]_{(x)}/(x^3)$.
Otrzymujemy bazę $1,x,x^2$ i lokalną krotność $3$.
Dwa przypadkowe rysunki tych krzywych nie wykażą,
że w liczbie przecięć należy uwzględnić właśnie trzy.
Wynika to z algebry lokalnego pierścienia.
\end{przyklad}

Założeniem podstawowym tego rozdziału jest, że $k$ jest
ciałem algebraicznie domkniętym charakterystyki zero.
Tam, gdzie jest to potrzebne, $X$ będzie gładką,
quasi-rzutową rozmaitością algebraiczną nad $k$,
a $S$ gładką, spójną, rzutową powierzchnią nad $k$.
Część konstrukcji zachowuje sens dla schematów osobliwych,
lecz wówczas iloczyn cykli wymaga narzędzi bardziej ogólnych
niż zwykłe przecięcia podrozmaitości.

\section{Cykle algebraiczne i równoważność wymierna}
\label{sec:ga7-cycles}

\begin{definicja}[Cykle wymiaru $r$]
\label{def:ga7-cycles}
Niech $X$ będzie schematem skończonego typu nad $k$, a $r\ge0$.
\emph{Cyklem algebraicznym wymiaru $r$} na $X$ nazywamy formalną
sumę $\alpha=\sum_{i=1}^m n_i[V_i]$, gdzie $n_i\in\Z$,
zaś $V_i\subseteq X$ są całkowitymi domkniętymi podschematami
wymiaru $r$ (każdy występuje co najwyżej raz).
Grupę takich cykli oznaczamy $Z_r(X)$.
Przymiotnik \emph{całkowity} oznacza nierozkładalny i zredukowany.
Cyklem zerowym nazywamy element $Z_0(X)$.
\end{definicja}

\begin{uwaga}[Dlaczego współczynniki mogą być ujemne?]
Punkt i krzywa są obiektami geometrycznymi, ale chcemy także
odejmować cykle, tak jak odejmujemy dzielniki.
Cykle efektywne mają współczynniki nieujemne.
Przejście do grupy pozwala pisać relacje, takie jak
$[p]-[q]$, i sprowadzać problem geometryczny do algebry abelowej.
Dla powierzchni całkowitej $S$ grupa $Z_1(S)$
jest po prostu grupą dzielników Weila wprowadzoną w
definicji~\ref{def:ga6-divisor-group}.
\end{uwaga}

\begin{definicja}[Dzielnik funkcji wymiernej na krzywej]
\label{def:ga7-rational-divisor}
Niech $W\subseteq X$ będzie całkowitym, domkniętym
podschematem wymiaru $r+1$, a $\nu:\widetilde W\to W$
jego normalizacją. Dla $f\in k(W)^\times$ definiujemy
$r$-cykl $\operatorname{div}_W(f)$ przez
\[
\operatorname{div}_W(f):=\nu_*
\bigl(\operatorname{div}_{\widetilde W}(f)\bigr).
\]
Ponieważ $\widetilde W$ jest normalna, jej lokalne
pierścienie w punktach kowymiaru jeden są pierścieniami
dyskretnego wartościowania. Współczynniki dzielnika
$\operatorname{div}_{\widetilde W}(f)$ to odpowiednie rzędy
zer i biegunów. Dla skończonego morfizmu $\nu$ obraz
składnika $[\widetilde V]$ ma współczynnik
$[\kappa(\widetilde V):\kappa(\nu(\widetilde V))]$ przy
$[\nu(\widetilde V)]$, gdzie $\widetilde V$ jest
całkowitym podschematem kowymiaru jeden.
W szczególności dla całkowitej krzywej $W$ mamy
\[
\operatorname{div}_W(f)=
\sum_{q\in\widetilde W_{\mathrm{cl}}}
[\kappa(q):\kappa(\nu(q))]\operatorname{ord}_q(f)[\nu(q)].
\]
W rozpatrywanej sytuacji nad algebraicznie domkniętym
$k$ stopnie ciał reszt są równe jeden. Jeśli jednak
$W$ jest osobliwa, nie należy przypisywać jej punktom
osobliwym rzędów wartościowania bez przejścia
do normalizacji; kilka punktów normalizacji może
mieć ten sam obraz w $W$.
\end{definicja}

\begin{przyklad}[Dwa punkty normalizacji nad jednym węzłem]
\label{prz:ga7-nodal-divisor}
Niech $k$ będzie ciałem algebraicznie domkniętym
charakterystyki zero i
$W=\V(y^2-x^2(x+1))\subseteq\A_k^2$.
Jest to krzywa całkowita z węzłem $p=(0,0)$.
Jej normalizację opisuje morfizm
$\nu:\A_k^1\to W$ zadany przez
$x=t^2-1$ oraz $y=t(t^2-1)$.
Równość $t=y/x$ w ciele funkcji pokazuje,
że morfizm jest biwymierny, a równanie
$t^2=x+1$ — że jest skończony.
Nad węzłem $p$ leżą dwa punkty $t=1$ i $t=-1$.
Ponieważ funkcja $x=t^2-1=(t-1)(t+1)$
ma na normalizacji dwa proste zera,
otrzymujemy $\operatorname{div}_W(x)=2[p]$.
Ten współczynnik pochodzi z sumowania rzędów
w dwóch punktach normalizacji, a nie z pojedynczego
wartościowania lokalnego pierścienia węzła.
\end{przyklad}

\begin{definicja}[Grupa Chowa]
\label{def:ga7-chow}
Niech $R_r(X)$ będzie podgrupą $Z_r(X)$ generowaną przez
wszystkie cykle $\operatorname{div}_W(f)$, gdzie
$W\subseteq X$ jest całkowitym domkniętym podschematem
wymiaru $r+1$ i $f\in k(W)^\times$.
\emph{Grupą Chowa cykli wymiaru $r$} nazywamy
$\operatorname{CH}_r(X)=Z_r(X)/R_r(X)$.
Dwa cykle reprezentujące ten sam element $\operatorname{CH}_r(X)$
nazywamy \emph{wymiernie równoważnymi}.
Jeżeli $X$ jest czystego wymiaru $n$, to piszemy
$\operatorname{CH}^i(X)=\operatorname{CH}_{n-i}(X)$.
Dla $r<0$ przyjmujemy $\operatorname{CH}_r(X)=0$.
\end{definicja}

\begin{przyklad}[Punkty na prostej rzutowej]
\label{prz:ga7-p1-points}
Na $\PP^1_k$ wprowadźmy współrzędną wymierną $t$ na
mapie afinicznej, a punkt w nieskończoności oznaczmy
$\infty$. Funkcja $t-a$ ma proste zero w punkcie $a\in k$
i prosty biegun w $\infty$, zatem
$\operatorname{div}(t-a)=[a]-[\infty]$.
W $\operatorname{CH}_0(\PP^1)$ wszystkie punkty zamknięte
mają tę samą klasę. Z drugiej strony każda funkcja wymierna
na gładkiej krzywej właściwej ma tyle samo zer co biegunów,
licząc z krotnościami. Stąd stopień
$\deg(\sum n_i[p_i])=\sum n_i$ nie zmienia się przy
równoważności wymiernej. Każdy cykl zerowy jest wymiernie
równoważny $d[\infty]$, a jego stopień wynosi $d$;
otrzymujemy izomorfizm $\operatorname{CH}_0(\PP^1)\cong\Z$.
\end{przyklad}

\begin{przyklad}[Dlaczego zwartość geometryczna ma znaczenie?]
\label{prz:ga7-a1-chow0}
Na $\A^1_k$ funkcja $t-a$ ma dzielnik $[a]$:
nie ma tu punktu w nieskończoności, w którym należałoby
zapisać biegun. Wobec tego $[a]=0$ w
$\operatorname{CH}_0(\A^1)$ i $\operatorname{CH}_0(\A^1)=0$.
Nie oznacza to, że punkt $a$ nie istnieje. Oznacza jedynie,
że grupa Chowa pamięta również, czy rozpatrujemy przestrzeń
właściwą, czy otwartą.
\end{przyklad}

\begin{stwierdzenie}[Stopień cyklu zerowego na przestrzeni właściwej]
\label{st:ga7-degree}
Niech $X$ będzie właściwym schematem skończonego typu nad
algebraicznie domkniętym ciałem $k$. Wtedy przyporządkowanie
$\deg(\sum n_p[p])=\sum n_p$ określa homomorfizm
$\deg:\operatorname{CH}_0(X)\to\Z$.
\end{stwierdzenie}
\begin{proof}
Trzeba sprawdzić, że stopień znika na generatorach
$R_0(X)$. Taki generator jest dzielnikiem funkcji
wymiernej na całkowitej krzywej właściwej $W\subseteq X$.
Przejdźmy do normalizacji $\widetilde W$.
Jest ona gładką właściwą krzywą, a $f$ wyznacza
odwzorowanie wymierne $\widetilde W\to\PP^1$, które
na gładkiej właściwej krzywej przedłuża się do morfizmu.
Dla funkcji niestałej stopień włókna nad $0$
i nad $\infty$ jest ten sam, równy stopniowi morfizmu;
stąd $\deg\operatorname{div}(f)=0$. Dla funkcji stałej
dzielnik jest zerowy. Ponieważ przedłużenie morfizmu
na właściwej gładkiej krzywej oraz równość stopni włókien
zostały już wykorzystane przy dzielnikach krzywych,
nie potrzeba tu nowego założenia.
\end{proof}

\begin{uwaga}[Gdzie ukryta jest trudność?]
O ile definicja $\operatorname{CH}_r(X)$ wymaga tylko
cykli i dzielników funkcji wymiernych, o tyle iloczyn dwóch
klas cykli nie jest prostą operacją na formalnych sumach:
podrozmaitości mogą przecinać się w zbyt dużym wymiarze.
Dla gładkich przestrzeni algebraicznych problem rozwiązuje
\emph{lemat o przesuwaniu cykli} i konstrukcja przecięcia
z uwzględnieniem krotności. Wyniki te jawnie wskażemy.
\end{uwaga}

\section{Od cykli do przecięć: morfizmy i iloczyn}
\label{sec:ga7-intersection}

\begin{definicja}[Przecięcie właściwe]
\label{def:ga7-proper-intersection}
Niech $X$ będzie gładką całkowitą rozmaitością wymiaru $n$,
a $V,W\subseteq X$ całkowitymi domkniętymi podrozmaitościami
o kowymiarach $a,b$. Mówimy, że $V,W$ \emph{przecinają się
właściwie}, gdy każda składowa nierozkładalna $V\cap W$
ma kowymiar $a+b$ w $X$; gdy $a+b>n$, oznacza to, że
przecięcie jest puste.
Nie należy mylić tego pojęcia z \emph{właściwością morfizmu}.
\end{definicja}

\begin{przyklad}[Wspólna składowa a brak przecięcia właściwego]
Na $\PP^2$ krzywe $\V_+(XY)$ i $\V_+(XZ)$
mają wspólną prostą $X=0$. Ich przecięcie zbiorowe
nie jest zerowymiarowe, choć stopnie obu równań wynoszą dwa.
Nie wolno więc sumować lokalnych długości w punktach,
aby dostać cztery: powstałaby suma po nieskończonej liczbie
punktów. W pierścieniu Chowa obie krzywe mają klasy $2h$,
a iloczyn ich klas to $4h^2$, gdzie $h$ jest klasą prostej.
Dostęp do takiego iloczynu wymaga właśnie równoważności
wymiernej i zastąpienia dzielników odpowiednimi przedstawicielami.
\end{przyklad}

\begin{twierdzenie}[Iloczyn na gładkiej rozmaitości — wynik konstrukcyjny]
\label{tw:ga7-chow-ring}
Niech $X$ będzie gładką quasi-rzutową rozmaitością
wymiaru $n$ nad ciałem $k$. Wówczas grupy Chowa mają kanoniczną, przemienną,
łączną i stopniowaną strukturę pierścienia z jedynką na
$\operatorname{CH}^*(X)=\bigoplus_{i=0}^n\operatorname{CH}^i(X)$,
w której $\alpha\cdot\beta\in\operatorname{CH}^{i+j}(X)$
dla $\alpha\in\operatorname{CH}^i(X)$,
$\beta\in\operatorname{CH}^j(X)$, a $1=[X]$.
Jeżeli całkowite podrozmaitości $V,W$ przecinają się
właściwie, ich iloczyn jest sumą klas składowych
przecięcia z odpowiednimi dodatnimi krotnościami
lokalnymi. Dla dwóch dzielników na gładkiej powierzchni,
które nie mają wspólnej składowej, są to długości
$\operatorname{length}(\mathcal O_{X,p}/(f,g))$ z
rozdziału~\ref{ch:ga6-powierzchnie}.
\end{twierdzenie}

\begin{uwaga}[Dokładny status dowodu]
\label{uw:ga7-moving-lemma}
Ogólne twierdzenie~\ref{tw:ga7-chow-ring} wykorzystuje
lemat o przesuwaniu cykli, teorię przecięć \emph{Gysina}
i dowód niezależności od wyboru przedstawicieli.
Samo wymaganie zgodności z przecięciami właściwymi nie jest
tutaj traktowane jako niezależny dowód istnienia i jednoznaczności
całej struktury pierścienia; chodzi o kanoniczny iloczyn
konstruowany w tej teorii.
Nie ma krótkiego dowodu wyłącznie z aksjomatów pierścienia.
Traktujemy jego istnienie i własności funktorialne jako
\emph{zewnętrzny wynik konstrukcyjny}. Natomiast wzory dla
dzielników na powierzchniach poniżej zostaną wyprowadzone
ze stopni wiązek na krzywych. Dzięki temu
rachunki z rozdziału 6 nie opierają się na niewyjaśnionej
operacji mnożenia symboli.
\end{uwaga}

\begin{stwierdzenie}[Przecięcie z dzielnikiem przez stopień]
\label{st:ga7-divisor-degree}
Niech $k$ będzie ciałem algebraicznie domkniętym
charakterystyki zero, $S$ gładką rzutową powierzchnią
nad $k$, $C\subseteq S$ całkowitą krzywą,
$D$ dzielnikiem na $S$, a $\nu:\widetilde C\to C$
normalizacją. Wtedy liczba przecięcia $D\cdot C$ jest równa
$\deg\bigl(\nu^*(\mathcal O_S(D)|_C)\bigr)$.
Jest addytywna względem $D$, nie zmienia się po
zastąpieniu $D$ dzielnikiem liniowo równoważnym i
dla $D=C$ i $C$ gładkiej daje $C^2=\deg N_{C/S}$.
\end{stwierdzenie}
\begin{proof}
Równość z definicji ogólnego iloczynu stanowi
zastosowanie pierwszej klasy Cherna do cyklu $[C]$;
wyprowadzimy jej definicję w sekcji~\ref{sec:ga7-chern}.
Wyjaśnijmy już teraz, dlaczego pozostałe własności są
konieczne. Dla dzielników $D,D'$ mamy
$\mathcal O_S(D+D')\cong\mathcal O_S(D)\otimes\mathcal O_S(D')$.
Ponieważ stopień iloczynu tensorowego wiązek na gładkiej
właściwej krzywej jest sumą stopni, otrzymujemy
$(D+D')\cdot C=D\cdot C+D'\cdot C$.
Jeżeli $D-D'=\operatorname{div}(f)$, wiązki
$\mathcal O_S(D),\mathcal O_S(D')$ są izomorficzne,
więc ich stopnie na $\widetilde C$ są równe.
W przypadku $C$ gładkiej z definicji normalnej
$N_{C/S}\cong\mathcal O_S(C)|_C$,
a więc $C^2=\deg N_{C/S}$.
Identyfikacja \emph{tego} stopnia z sumą lokalnych
krotności przecięć dla różnych krzywych jest
istotnym twierdzeniem porównawczym teorii przecięć;
nie wynika z samej addytywności stopnia.
\end{proof}

\begin{przyklad}[Długości zamiast odgadywania krotności]
Niech $m\ge1$, $C=\V(y)$ i $D=\V(y-x^m)$ w $\A^2_k$.
W początku oba równania generują ideał $(y,x^m)$.
Domknięcia rzutowe tych krzywych mają równania
$Y=0$ oraz $YZ^{m-1}-X^m=0$, stopni odpowiednio
$1$ i $m$. Twierdzenie Bézouta daje łączną liczbę
przecięć równą $m$.
Sprawdźmy, że w nieskończoności nie ma przecięć.
Dla $m\ge2$ na $Z=0$ drugie równanie daje $X=0$,
zatem jedynym punktem drugiej krzywej w nieskończoności
jest $[0:1:0]$, który nie należy do prostej $Y=0$.
Dla $m=1$ drugie równanie ma postać $Y-X=0$
i jego punkt w nieskończoności to $[1:1:0]$;
on również nie należy do prostej $Y=0$.
Zatem cały wkład do liczby przecięć pochodzi
z początku $(0,0)$ i jest równy $m$.
\end{przyklad}

\begin{center}
\begin{minipage}{\linewidth}
\centering
\begin{tikzpicture}[>=Stealth,every node/.style={font=\small}]
 \node[draw,rounded corners,text width=3.0cm,align=center,minimum height=1.35cm] (a) at (-3.65,0)
 {Dwie krzywe bez wspólnej składowej\\$C\cap D$ skończony};
 \node[draw,rounded corners,text width=3.0cm,align=center,minimum height=1.35cm] (b) at (3.65,0)
 {Suma długości lokalnych\\$\sum_p i_p(C,D)$};
 \draw[->,thick] (a.east)--node[above,font=\footnotesize] {rachunek lokalny} (b.west);
 \node[draw,rounded corners,text width=5.4cm,align=center,minimum height=0.95cm] (c) at (0,-2.0)
 {Globalna klasa $[C]\cdot[D]$ w $\operatorname{CH}^2(S)$};
 \draw[->,thick] (a.south)--(c.north west);
 \draw[->,thick] (b.south)--(c.north east);
\end{tikzpicture}

{\small Rysunek 7.1. Dwie drogi do liczby przecięcia: lokalne pierścienie i globalne klasy cykli.}
\end{minipage}
\end{center}

\section{Podstawowe grupy Chowa: przestrzeń rzutowa}
\label{sec:ga7-projective-chow}

\begin{twierdzenie}[Pierścień Chowa przestrzeni rzutowej]
\label{tw:ga7-chow-pn}
Niech $k$ będzie ciałem algebraicznie domkniętym,
a $n\ge0$ liczbą całkowitą. Dla $n\ge1$ oznaczmy
przez $h=[\PP^{n-1}]\in\operatorname{CH}^1(\PP^n_k)$
klasę hiperpłaszczyzny. Wówczas
\[
\operatorname{CH}^*(\PP^n_k)\cong\Z[h]/(h^{n+1}).
\]
Jeżeli $n=0$, lewą stroną jest $\Z$, a zapis po prawej
oznacza $\Z[h]/(h)\cong\Z$, z formalnym $h=0$.
Dla $n\ge1$ klasa $h^j$ jest klasą liniowej podprzestrzeni
kowymiaru $j$; we wszystkich przypadkach $\deg(h^n)=1$
(przy $n=0$ jest to stopień $[\PP^0_k]$).
\end{twierdzenie}
\begin{proof}
Dla $n=0$ mamy $\PP^0_k=\Spec k$ oraz
$\operatorname{CH}^0(\PP^0_k)=\Z[\PP^0_k]$,
więc twierdzenie zachodzi z $h=0$.
Załóżmy dalej, że $n\ge1$.
Przyjmujemy dwa standardowe fakty o grupach Chowa:
\emph{ciąg lokalizacji} dla domkniętego zanurzenia
$Z\hookrightarrow X$ z otwartym dopełnieniem $U$,
$\operatorname{CH}_r(Z)\to\operatorname{CH}_r(X)
\to\operatorname{CH}_r(U)\to0$, oraz
\emph{niezmienniczość homotopijną}, z której wynika
$\operatorname{CH}_r(\A^n)=0$ dla $r<n$ i
$\operatorname{CH}_n(\A^n)=\Z[\A^n]$.
Dla $X=\PP^n$, $Z=\PP^{n-1}$ będącej hiperpłaszczyzną,
$U=\A^n$ dostajemy, dla $r<n$, surjekcję
$\operatorname{CH}_r(\PP^{n-1})\to\operatorname{CH}_r(\PP^n)$.
Indukcja po $n$ pokazuje, że $\operatorname{CH}_r(\PP^n)$
jest generowana przez klasę liniowej $\PP^r$.
Dla $r=n$ generatorem jest $[\PP^n]$.
Nie może być niezerowego stosunku $m[\PP^r]=0$:
po przecięciu z $r$ ogólnymi hiperpłaszczyznami otrzymujemy
$m$ razy klasę punktu, o stopniu $m$.
Homomorfizm stopnia z
stwierdzenia~\ref{st:ga7-degree} pokazuje, że $m=0$.
Zatem wszystkie te grupy są izomorficzne z $\Z$.
Iloczyn $j$ ogólnych hiperpłaszczyzn jest liniową
podprzestrzenią kowymiaru $j$, zatem jej klasą jest $h^j$.
Przy $j=n+1$ wymiar cyklu stałby się ujemny, więc
$h^{n+1}=0$. Nie ma innych relacji, gdyż
$1,h,\ldots,h^n$ generują niezależne grupy
w różnych stopniach.
\end{proof}

\begin{przyklad}[Bézout jako mnożenie klas]
\label{prz:ga7-bezout}
Niech $d\ge1$, a $F\ne0$ będzie jednorodnym
wielomianem stopnia $d$. Wówczas $F$ definiuje
w $\PP^n$ efektywny dzielnik klasy $dh$.
Jeżeli $F_1,\ldots,F_n$ są niezerowymi jednorodnymi
wielomianami dodatnich stopni $d_1,\ldots,d_n$
i ich wspólny zbiór zer jest zerowymiarowy (bez składowych
wymiaru dodatniego), to jego cykl przecięcia ma klasę
$(d_1\cdots d_n)h^n$. Stopień tego cyklu jest
$d_1\cdots d_n$. Dla dwóch krzywych płaskich daje to
$de$, a dla trzech powierzchni w $\PP^3$ --- $d_1d_2d_3$.
Poszczególne punkty przecięcia mogą mieć różne krotności,
lecz ich suma jest ustalona.
\end{przyklad}

\begin{uwaga}[Co zostało dowiedzione samodzielnie?]
Po przyjęciu ciągu lokalizacji, niezmienniczości
homotopijnej i konstrukcji iloczynu dla gładkich
rozmaitości dowód pierścienia Chowa $\PP^n$ jest
już rachunkiem elementarnym. Nie twierdzimy, że
owe trzy twierdzenia funktorialne zostały dowiedzione
od podstaw w tym skrypcie.
\end{uwaga}

\section{Klasy Cherna: algebraiczny opis skręcenia wiązki}
\label{sec:ga7-chern}

Wiązki liniowe i ich dzielniki mierzyły już, ile zer
ma przekrój na krzywej. Dla wiązki wektorowej rzędu
większego niż jeden pojedynczy dzielnik nie wystarcza.
Klasy Cherna są wielowymiarowym uogólnieniem stopnia:
przypisują wiązce $\mathcal E$ kolejne klasy
$c_i(\mathcal E)$ w pierścieniu Chowa.

\begin{definicja}[Pierwsza klasa Cherna]
\label{def:ga7-c1}
Niech $X$ będzie gładką quasi-rzutową rozmaitością,
a $L$ wiązką liniową na $X$. Jeśli $L\cong\mathcal O_X(D)$
dla dzielnika Cartiera $D$, definiujemy
$c_1(L)=[D]\in\operatorname{CH}^1(X)$.
Jest to poprawna definicja, ponieważ dzielniki liniowo
równoważne reprezentują tę samą klasę, a na gładkiej
rozmaitości $\operatorname{Pic}(X)\cong\operatorname{CH}^1(X)$.
Dla wiązki trywialnej $c_1(\mathcal O_X)=0$.
\end{definicja}

\begin{stwierdzenie}[Pierwsza klasa Cherna i iloczyn tensorowy]
\label{st:ga7-c1-add}
Niech $X$ będzie gładką quasi-rzutową rozmaitością,
a $L,M$ wiązkami liniowymi na $X$. Wtedy
$c_1(L\otimes M)=c_1(L)+c_1(M)$ oraz
$c_1(L^\vee)=-c_1(L)$.
\end{stwierdzenie}
\begin{proof}
Wybierzmy dzielniki $D,E$ z $L\cong\mathcal O_X(D)$
i $M\cong\mathcal O_X(E)$. Z konstrukcji snopa dzielnika
wynika $L\otimes M\cong\mathcal O_X(D+E)$, a
$L^\vee\cong\mathcal O_X(-D)$.
Przejście do klas cykli daje obie równości.
\end{proof}

\begin{twierdzenie}[Formuła wiązki rzutowej — wynik zewnętrzny]
\label{tw:ga7-projective-bundle}
Niech $X$ będzie gładką quasi-rzutową rozmaitością,
$\mathcal E$ wiązką wektorową rzędu $r\ge1$, a
$\pi:\PP_X(\mathcal E)\to X$ przestrzenią
jednowymiarowych ilorazów włókien $\mathcal E$.
Oznaczmy $\xi=c_1(\mathcal O_{\PP(\mathcal E)}(1))$.
Wówczas $\operatorname{CH}^*(\PP_X(\mathcal E))$ jest
wolnym modułem nad $\operatorname{CH}^*(X)$ o bazie
$1,\xi,\ldots,\xi^{r-1}$; działanie współczynników
jest dane przez $\pi^*$.
\end{twierdzenie}

\begin{definicja}[Wyższe klasy Cherna]
\label{def:ga7-c-classes}
Niech $X$ będzie gładką quasi-rzutową rozmaitością nad $k$,
$\mathcal E$ wiązką wektorową rzędu $r\ge1$,
$\pi:\PP_X(\mathcal E)\to X$ projektywizacją
jednowymiarowych ilorazów, a
$\xi=c_1(\mathcal O_{\PP(\mathcal E)}(1))$.
Z formuły wiązki rzutowej~\ref{tw:ga7-projective-bundle}
istnieją jednoznaczne klasy
$c_i(\mathcal E)\in\operatorname{CH}^i(X)$, $1\le i\le r$,
takie, że
\[
\xi^r-\pi^*c_1(\mathcal E)\xi^{r-1}
+\pi^*c_2(\mathcal E)\xi^{r-2}
-\cdots+(-1)^r\pi^*c_r(\mathcal E)=0.
\]
Definiujemy także $c_0(\mathcal E)=1$ i
$c_i(\mathcal E)=0$ dla $i>r$.
\emph{Całkowitą klasą Cherna} nazywamy
$c(\mathcal E)=\sum_{i=0}^r c_i(\mathcal E)$.
\end{definicja}

\begin{uwaga}[Skąd znaki w definicji?]
Przy naszej konwencji ilorazów, jeśli $\mathcal E=L$
ma rząd jeden, to $\PP_X(L)=X$ i
$\mathcal O_{\PP(L)}(1)\cong L$. Równanie staje się
$\xi-c_1(L)=0$, zgodnie z wcześniejszą definicją.
Dla sumy dwóch wiązek liniowych $L\oplus M$ relacja
przyjmuje postać
$\xi^2-(c_1(L)+c_1(M))\xi+c_1(L)c_1(M)=0$.
Różne konwencje projektywizacji zmieniają postać
równań; tutaj będziemy konsekwentnie używać ilorazów.
\end{uwaga}

\begin{twierdzenie}[Własności klas Cherna]
\label{tw:ga7-chern-axioms}
Niech $X$ będzie gładką quasi-rzutową rozmaitością.
Dla wiązek wektorowych $\mathcal E,\mathcal F$ na $X$
zachodzą następujące własności:
\begin{enumerate}[label=(\roman*)]
\item $c_1(\mathcal E)=c_1(\det\mathcal E)$.
\item Jeśli $0\to\mathcal E'\to\mathcal E\to\mathcal E''\to0$
      jest dokładnym ciągiem wiązek wektorowych, to
      $c(\mathcal E)=c(\mathcal E')c(\mathcal E'')$
      (\emph{wzór Whitneya}).
\item Dla morfizmu $f:Y\to X$ między gładkimi
      quasi-rzutowymi rozmaitościami zachodzi
      $c_i(f^*\mathcal E)=f^*c_i(\mathcal E)$.
\item Jeżeli $\mathcal E=L_1\oplus\cdots\oplus L_r$
      jest sumą wiązek liniowych, to
      $c(\mathcal E)=\prod_{j=1}^r(1+c_1(L_j))$.
\end{enumerate}
\end{twierdzenie}
\begin{proof}[Dowód z wykorzystaniem zasady rozszczepienia]
Najpierw wyjaśnimy zasadę rozszczepienia. Dla wiązki
$\mathcal E$ można skonstruować kolejne projektywizacje
wiązek, prowadzące do morfizmu $p:Y\to X$, dla którego
$p^*:\operatorname{CH}^*(X)\to\operatorname{CH}^*(Y)$
jest iniektywne, a $p^*\mathcal E$ posiada filtrację
przez podwiązki, której ilorazy są wiązkami liniowymi
$L_1,\ldots,L_r$. Iniektywność na każdym etapie
wynika z formuły wiązki rzutowej: pullback
identyfikuje pierścień z podmodułem współczynnika
przy $\xi^0$ w swobodnej bazie $1,\xi,\ldots,\xi^{r-1}$.
Pełna konstrukcja flag przez kolejne uniwersalne ilorazy
stanowi standardową część teorii wiązek rzutowych.

Dla sumy wiązek liniowych definicja przez relację
projektową daje
$c(\mathcal E)=\prod_j(1+x_j)$,
gdzie $x_j=c_1(L_j)$; dla filtracji z ilorazami liniowymi
obowiązuje ten sam wzór. Jego wykazanie to podstawowy
etap konstrukcyjny teorii klas Cherna, który
przyjmujemy wraz z formułą wiązki rzutowej.
Należy podkreślić, że sama iniektywność odwzorowania indukowanego
nie dowodzi jeszcze tej tożsamości: jej uzasadnienie
wymaga zgodności klas Cherna z filtracjami wiązek.
Wyrażenia $c_i$ stają się więc elementarnymi funkcjami
symetrycznymi zmiennych $x_j$, zwanych
\emph{formalnymi pierwiastkami Cherna}.

Jeżeli $0\to\mathcal E'\to\mathcal E\to\mathcal E''\to0$
jest dokładny, wybierzmy wspólną przestrzeń rozszczepienia
dla $\mathcal E'$ i $\mathcal E''$, z iniektywnym
odwzorowaniem indukowanym na grupach Chowa. Filtracja
$\mathcal E'$ i przeciwobrazy podwiązek z filtracji
$\mathcal E''$ dają filtrację $\mathcal E$,
której ilorazy są łącznie ilorazami obu wiązek.
W tej przestrzeni mamy zatem
$c(\mathcal E)=c(\mathcal E')c(\mathcal E'')$;
iniektywność odwzorowania indukowanego daje (ii) na $X$.
Przypadek rozszczepionego ciągu daje (iv).
Ponieważ wyznacznik ma pierwszą klasę Cherna równą
sumie pierwszych klas ilorazów, otrzymujemy również (i).

Własność (iii) można odczytać bezpośrednio z definicji:
projektywizacja i tautologiczna wiązka ilorazowa
są zgodne ze zmianą bazy, a pullback zachowuje relację
wielomianową definiującą klasy $c_i$. Jednoznaczność
współczynników zapewniona przez formułę wiązki rzutowej
daje $c_i(f^*\mathcal E)=f^*c_i(\mathcal E)$.
W ten sposób dowód redukuje się do \emph{dwóch} jawnych
wyników konstrukcyjnych: formuły wiązki rzutowej
i zasady rozszczepienia, zamiast przyjmować osobno
każdą z czterech własności.
\end{proof}

\begin{przyklad}[Wiązka sumy dwóch liniowych na $\PP^2$]
\label{prz:ga7-chern-sum}
Niech $h=c_1(\mathcal O_{\PP^2}(1))$,
$\mathcal E=\mathcal O_{\PP^2}(2)\oplus\mathcal O_{\PP^2}(3)$.
Ponieważ $c(\mathcal O(d))=1+dh$, wzór Whitneya daje
$c(\mathcal E)=(1+2h)(1+3h)=1+5h+6h^2$.
Zatem $c_1(\mathcal E)=5h$ i $c_2(\mathcal E)=6h^2$.
Liczba $\deg c_2(\mathcal E)=6$ ma również
interpretację jako liczba zer ogólnego przekroju
tej wiązki, liczonych z krotnościami (pod warunkiem,
że zbiór zer ma oczekiwany wymiar zero).
\end{przyklad}

\begin{stwierdzenie}[Zera przekroju i najwyższa klasa Cherna]
\label{st:ga7-zero-section}
Niech $X$ będzie gładką rzutową rozmaitością wymiaru $n$,
$\mathcal E$ wiązką wektorową rzędu $n$ na $X$,
a $s$ jej regularnym przekrojem, którego schemat zer
$Z(s)$ jest zerowymiarowy. Wtedy
$[Z(s)]=c_n(\mathcal E)\cap[X]$ w $\operatorname{CH}_0(X)$,
a zatem długość $Z(s)$ wynosi $\deg c_n(\mathcal E)$.
\end{stwierdzenie}
\begin{uwaga}[Dlaczego warunek regularności jest potrzebny?]
Przekrój może zerować się na całej krzywej albo
być identycznie zerowy. Wówczas nie można po prostu
liczyć jego zer jako izolowanych punktów. Twierdzenie
o klasie najwyższej wymaga przekroju regularnego,
tj. takiego, którego schemat zer ma kowymiar równy
rzędowi wiązki (lub pustego zbioru zer, gdy klasa znika).
Jego dowód w pełnej ogólności korzysta z teorii
przecięć regularnych zanurzeń; oznaczamy go jako
\emph{zewnętrzny}. Konkretny przykład poniżej
możemy jednak obliczyć bezpośrednio.
\end{uwaga}

\section{Wiązki styczne przestrzeni rzutowych}
\label{sec:ga7-tangent}

\begin{twierdzenie}[Ciąg Eulera i klasy Cherna $\PP^n$]
\label{tw:ga7-euler}
Niech $k$ będzie ciałem algebraicznie domkniętym
charakterystyki zero, $n\ge1$ i
$h=c_1(\mathcal O_{\PP^n_k}(1))$.
Wówczas istnieje dokładny ciąg wiązek
\[
0\longrightarrow\mathcal O_{\PP^n}
\longrightarrow\mathcal O_{\PP^n}(1)^{\oplus(n+1)}
\longrightarrow T_{\PP^n}\longrightarrow0,
\]
a całkowita klasa Cherna wiązki stycznej jest równa
$c(T_{\PP^n})=(1+h)^{n+1}$ w
$\Z[h]/(h^{n+1})$. W szczególności
$c_i(T_{\PP^n})=\binom{n+1}{i}h^i$ dla $0\le i\le n$.
\end{twierdzenie}
\begin{proof}
Oznaczmy przez $V=k^{n+1}$ przestrzeń, której proste
parametryzuje $\PP^n=\PP(V)$ w konwencji prostych.
Wybór niezerowego wektora $v\in\ell$ reprezentującego
$[\ell]$ daje opis
$T_{[\ell]}\PP(V)\cong\operatorname{Hom}(\ell,V/\ell)$:
deformacja prostej jest odwzorowaniem z jej kierunku
do kierunków poprzecznych. Na całej $\PP(V)$ otrzymujemy
ciąg wiązek
$0\to\mathcal O\to\mathcal O(1)\otimes V
\to T_{\PP(V)}\to0$, przy czym pierwszy morfizm
wysyła lokalnie stałą $1$ na sekcję złożoną z
współrzędnych jednorodnych. Jego obraz jest podwiązką,
a iloraz ma włókno $\operatorname{Hom}(\ell,V/\ell)$.
To daje ciąg Eulera. Ze wzoru Whitneya,
$c(\mathcal O)=1$ i $c(\mathcal O(1))=1+h$
otrzymujemy $c(T_{\PP^n})=(1+h)^{n+1}$.
Ponieważ $h^{n+1}=0$, rozwinięcie dwumianowe
kończy się w stopniu $n$.
\end{proof}

\begin{przyklad}[Trzy zera pola wektorowego na $\PP^2$]
\label{prz:ga7-vector-field}
Z twierdzenia~\ref{tw:ga7-euler} mamy
$c_2(T_{\PP^2})=3h^2$, a więc jego stopień wynosi $3$.
Rozważmy algebraiczne działanie grupy multiplikatywnej
$\mathbb G_m$ na $\PP^2_k$ dane przez
$\lambda\cdot[x:y:z]=[x:\lambda y:\lambda^2 z]$.
Różniczkując to działanie w jedynce $\lambda=1$,
otrzymujemy globalne algebraiczne pole wektorowe.
Nad $k=\C$ tę samą konstrukcję można zapisać
analitycznie, podstawiając $\lambda=e^t$
i różniczkując względem $t$ w zerze.
Pole znika dokładnie w trzech punktach współrzędnych
$[1:0:0]$, $[0:1:0]$, $[0:0:1]$.
W pobliżu pierwszego punktu na mapie $x\ne0$
współrzędne $(u=y/x,v=z/x)$ spełniają
$\dot u=u$, $\dot v=2v$; zatem zero jest proste.
Na pozostałych dwóch mapach otrzymujemy odpowiednio
wagi $(-1,1)$ oraz $(-2,-1)$ i znowu proste zera.
Łączna długość schematu zer wynosi $1+1+1=3$,
co dokładnie zgadza się z $\deg c_2(T_{\PP^2})$.
\end{przyklad}

\begin{center}
\begin{minipage}{\linewidth}
\centering
\begin{tikzpicture}[>=Stealth,scale=0.95,every node/.style={font=\small}]
 \coordinate (A) at (-2.8,0); \coordinate (B) at (2.8,0); \coordinate (C) at (0,2.95);
 \draw[thick] (A)--(B)--(C)--cycle;
 \fill[blue!70!black] (A) circle (2.6pt);
 \fill[blue!70!black] (B) circle (2.6pt);
 \fill[blue!70!black] (C) circle (2.6pt);
 \node[below left=4pt of A,font=\footnotesize] {$[1:0:0]$};
 \node[below right=4pt of B,font=\footnotesize] {$[0:1:0]$};
 \node[above=7pt of C,font=\footnotesize] {$[0:0:1]$};
 \node[align=center,font=\footnotesize] at (0,0.87) {trzy zera};
\end{tikzpicture}

{\small Rysunek 7.2. Schematyczne położenie trzech zer algebraicznego pola wektorowego na $\PP^2$; trójkąt przedstawia jedynie współrzędne wierzchołkowe, nie całą rzeczywistą geometrię $\PP^2(\C)$.}
\end{minipage}
\end{center}

\section{Gładkie hiperpowierzchnie i ich niezmienniki}
\label{sec:ga7-hypersurfaces}

\begin{twierdzenie}[Klasy Cherna hiperpowierzchni]
\label{tw:ga7-hypersurface-chern}
Niech $X=\V_+(F)\subseteq\PP^n_k$ będzie gładką
hiperpowierzchnią stopnia $d\ge1$, gdzie $n\ge2$.
Oznaczmy $h=c_1(\mathcal O_X(1))$.
W pierścieniu $\operatorname{CH}^*(X)$ zachodzi
\[
 c(T_X)=\frac{(1+h)^{n+1}}{1+dh},
\]
przy czym prawą stronę rozumiemy jako szereg formalny
obcięty w stopniu większym niż $\dim X=n-1$.
W szczególności $c_1(T_X)=(n+1-d)h$.
\end{twierdzenie}
\begin{proof}
Ponieważ $X$ jest gładka, jest regularną hiperpowierzchnią,
a wiązka normalna jej zanurzenia w $\PP^n$ jest
$N_{X/\PP^n}\cong\mathcal O_X(d)$; wynika to
z tego, że ideał $X$ jest lokalnie generowany przez $F$,
zaś jego snop to $\mathcal O_{\PP^n}(-d)$.
Mamy więc dokładny ciąg wiązek stycznych
$0\to T_X\to T_{\PP^n}|_X\to\mathcal O_X(d)\to0$.
Ze wzoru Whitneya otrzymujemy
$c(T_{\PP^n}|_X)=c(T_X)c(\mathcal O_X(d))$.
Naturalność klas Cherna i ciąg Eulera dają
$c(T_{\PP^n}|_X)=(1+h)^{n+1}$, natomiast
$c(\mathcal O_X(d))=1+dh$.
Ponieważ element $dh$ ma dodatni stopień, $1+dh$
ma odwrotność w skończenie stopniowanym pierścieniu
Chowa $X$. Dzieląc, otrzymujemy wzór.
\end{proof}

\begin{wniosek}[Gładkie powierzchnie stopnia $d$ w $\PP^3$]
\label{wn:ga7-degree-d-surface}
Niech $k$ będzie ciałem algebraicznie domkniętym
charakterystyki zero, $S\subseteq\PP^3_k$ gładką
powierzchnią stopnia $d\ge1$, $H$ dzielnikiem
hiperpłaszczyznowym na $S$ i $h=[H]=c_1(\mathcal O_S(1))$.
Dla dowolnego dzielnika kanonicznego $K_S$ zachodzą
\[
 K_S\sim(d-4)H,\qquad K_S^2=d(d-4)^2,\qquad
 c_2(T_S)=(d^2-4d+6)h^2,
\]
a $\deg c_2(T_S)=d(d^2-4d+6)$.
\end{wniosek}
\begin{proof}
Dla $n=3$ rozwijamy iloraz
$(1+h)^4/(1+dh)$ do drugiego stopnia:
$(1+4h+6h^2)(1-dh+d^2h^2)
=1+(4-d)h+(d^2-4d+6)h^2$.
Ponadto z dokładnego ciągu wiązek normalnych
$0\to T_S\to T_{\PP^3}|_S\to\mathcal O_S(d)\to0$
i tożsamości $\det T_{\PP^3}\cong\mathcal O_{\PP^3}(4)$
wynika po wzięciu wyznaczników
$\det T_S\cong\mathcal O_S(4-d)$, a więc
$\omega_S\cong\mathcal O_S(d-4)$.
Jest to izomorfizm wiązek liniowych, dlatego
$K_S\sim(d-4)H$, nie tylko równość ich klas w grupie Chowa.
Ponieważ $S$ ma stopień $d$, dwa ogólne przekroje
hiperpłaszczyzną przecinają się na $S$ w cyklu zerowym
stopnia $d$, tzn. $\deg(h^2)=d$.
Zatem $K_S^2=(d-4)^2\deg(h^2)=d(d-4)^2$,
a stopień drugiej klasy Cherna jest równy
$d(d^2-4d+6)$.
\end{proof}

\begin{przyklad}[Cztery powierzchnie, cztery różne drugie klasy]
\label{prz:ga7-surface-invariants}
Dla $d=1,2,3,4$ wzory dają odpowiednio:
\[
\begin{array}{c|c|c|c}
 d & \text{powierzchnia} & K_S^2 & \deg c_2(T_S)\\ \hline
 1 & \PP^2 & 9 & 3\\
 2 & \text{gładka kwadryka} & 8 & 4\\
 3 & \text{gładka powierzchnia kubiczna} & 3 & 9\\
 4 & \text{gładka kwartyka K3} & 0 & 24
\end{array}
\]
Wartość $24$ dla kwartyki K3 nie jest osobnym aksjomatem
definicji K3: wynika z obliczenia wiązki stycznej.
\end{przyklad}

\section{Riemann--Roch na powierzchniach: dowód z twierdzenia dla krzywych}
\label{sec:ga7-rr-surface}

Tutaj możemy udowodnić ważne twierdzenie, nie przyjmując
od razu całej teorii klas charakterystycznych. Punktem wyjścia
będzie Riemann--Roch dla \emph{gładkiej krzywej}, którego
dowodowe podstawy zostały omówione w rozdziale 4.

\begin{definicja}[Charakterystyka Eulera snopa koherentnego]
\label{def:ga7-chi}
Niech $S$ będzie gładką rzutową powierzchnią, a
$\mathcal F$ snopem koherentnym na $S$. Definiujemy
$\chi(S,\mathcal F)=\sum_{i=0}^{2}(-1)^i
\dim_k H^i(S,\mathcal F)$.
Wymiary są skończone dzięki twierdzeniu o skończoności
kohomologii koherentnej na schemacie właściwym.
Gdy $\mathcal F=\mathcal O_S(D)$, piszemy
$\chi(\mathcal O_S(D))$.
\end{definicja}

\begin{lemat}[Addytywność charakterystyki Eulera]
\label{lem:ga7-chi-exact}
Niech $0\to\mathcal F'\to\mathcal F\to\mathcal F''\to0$
będzie dokładnym ciągiem snopów koherentnych na
gładkiej rzutowej powierzchni $S$. Wtedy
$\chi(\mathcal F)=\chi(\mathcal F')+\chi(\mathcal F'')$.
\end{lemat}
\begin{proof}
Ciąg krótki indukuje długi ciąg kohomologii
$H^0(\mathcal F')\to H^0(\mathcal F)\to\cdots
\to H^2(\mathcal F'')\to0$.
Jest to skończony ciąg dokładny skończenie wymiarowych
przestrzeni wektorowych. Suma naprzemienna ich wymiarów
wynosi zero: jeśli $V_i$ są kolejnymi wyrazami, to
$\dim V_i=\dim\operatorname{im}(V_{i-1}\to V_i)+
\dim\operatorname{im}(V_i\to V_{i+1})$, a po zsumowaniu
naprzemiennym wszystkie wyrazy obrazu znoszą się parami.
Zgrupowanie wyrazów według trzech snopów daje tezę.
\end{proof}

\begin{lemat}[Przyrost charakterystyki po dodaniu krzywej]
\label{lem:ga7-rr-increment}
Niech $k$ będzie ciałem algebraicznie domkniętym charakterystyki
zero, $S$ gładką rzutową powierzchnią nad $k$,
$K_S$ dowolnym dzielnikiem kanonicznym na $S$,
$D$ dowolnym dzielnikiem, a $C\subseteq S$ gładką
całkowitą krzywą. Wtedy
\[
\chi(\mathcal O_S(D+C))-\chi(\mathcal O_S(D))
= D\cdot C+\frac{C^2-K_S\cdot C}{2}.
\]
\end{lemat}
\begin{proof}
Krzywa $C$ jest efektywnym dzielnikiem Cartiera.
Pomnożenie ciągu
$0\to\mathcal O_S(-C)\to\mathcal O_S\to\mathcal O_C\to0$
przez wiązkę liniową $\mathcal O_S(D+C)$ jest dokładne,
ponieważ iloczyn tensorowy z wiązką lokalnie wolną
zachowuje dokładność. Otrzymujemy
\[
0\longrightarrow\mathcal O_S(D)
\longrightarrow\mathcal O_S(D+C)
\longrightarrow\mathcal O_C(D+C)\longrightarrow0.
\]
Lemat~\ref{lem:ga7-chi-exact} daje
$\chi(\mathcal O_S(D+C))-\chi(\mathcal O_S(D))
=\chi(\mathcal O_C(D+C))$; tu korzystamy także z tego,
że kohomologie snopa $i_*L$ na $S$ są kohomologiami
$L$ na $C$, gdzie $i:C\hookrightarrow S$ jest zanurzeniem.
Riemann--Roch dla gładkiej właściwej krzywej
$C$ mówi, że dla wiązki $L$ na $C$ zachodzi
$\chi(C,L)=\deg L+1-g(C)$. Po podstawieniu
$L=\mathcal O_C(D+C)$ otrzymujemy
$\chi= D\cdot C+C^2+1-g(C)$.
Wzór adjunkcji z twierdzenia~\ref{tw:ga6-adjunction}
daje $1-g(C)=-\frac12(C^2+K_S\cdot C)$.
Łącząc te równości, uzyskujemy tezę.
\end{proof}

\begin{twierdzenie}[Riemann--Roch dla dzielników na powierzchni]
\label{tw:ga7-rr-surfaces}
Niech $S$ będzie gładką spójną rzutową powierzchnią
nad algebraicznie domkniętym ciałem charakterystyki zero,
$K_S$ jej dzielnikiem kanonicznym, a $D$ dowolnym dzielnikiem.
Wtedy
\[
\chi(\mathcal O_S(D))=\chi(\mathcal O_S)
+\frac{D\cdot(D-K_S)}{2}.
\]
\end{twierdzenie}
\begin{proof}
Zdefiniujmy funkcję na grupie dzielników
\[
Q(D)=\chi(\mathcal O_S(D))-\chi(\mathcal O_S)
-\frac{D^2-K_S\cdot D}{2}.
\]
Z lematu~\ref{lem:ga7-rr-increment} wynika, że
dla każdej gładkiej całkowitej krzywej $C$ mamy
$Q(D+C)=Q(D)$, ponieważ
\[
\frac{(D+C)^2-K_S\cdot(D+C)}{2}
-\frac{D^2-K_S\cdot D}{2}
=D\cdot C+\frac{C^2-K_S\cdot C}{2}.
\]
W szczególności, podstawiając $D-C$ za $D$,
dostajemy również $Q(D-C)=Q(D)$.

Pozostaje sprawdzić, że każdy dzielnik jest liniowo
równoważny różnicy dwóch gładkich całkowitych krzywych.
Wybierzmy bardzo obszerny dzielnik $H$.
Dla dostatecznie dużego $m$ zarówno $mH$,
jak i $D+mH$ są bardzo obszerne (jest to znany
wynik Serre'a o obszerności; odróżniamy go od
samego twierdzenia o znikaniu kohomologii).
Na mocy twierdzenia Bertiniego, nad ciałem
algebraicznie domkniętym charakterystyki zero
ogólny element każdego z tych liniowych układów
jest gładką całkowitą krzywą, oznaczmy je
$C_0\sim mH$, $C_1\sim D+mH$.
Stąd $D\sim C_1-C_0$.

Wielkość $Q$ zależy tylko od klasy liniowej,
ponieważ zarówno $\mathcal O_S(D)$, jak i liczby
przecięć nie zmieniają się przy takiej równoważności.
Dzięki dwukrotnemu zastosowaniu równości
$Q(D\pm C)=Q(D)$ mamy
$Q(D)=Q(C_1-C_0)=Q(0)=0$.
Otrzymujemy żądany wzór. Wykorzystane rezultaty
zewnętrzne to Riemann--Roch na krzywej,
Bertini i dostateczna obszerność $D+mH$;
nie używaliśmy Hirzebrucha--Riemanna--Rocha.
\end{proof}

\begin{wniosek}[Parzystość wyrażenia przecięciowego]
\label{wn:ga7-parity}
Niech $S$ będzie gładką rzutową powierzchnią
nad algebraicznie domkniętym ciałem charakterystyki zero,
a $K_S$ jej dowolnym dzielnikiem kanonicznym.
Dla każdego dzielnika $D$ na $S$ liczba
$D\cdot(D-K_S)$ jest parzysta.
\end{wniosek}
\begin{proof}
Z twierdzenia~\ref{tw:ga7-rr-surfaces} wynika,
że połowa tej liczby jest różnicą dwóch
charakterystyk Eulera, a więc liczbą całkowitą.
\end{proof}

\begin{przyklad}[Riemann--Roch na $\PP^2$ dla dowolnego $d\in\Z$]
\label{prz:ga7-rr-p2}
Na $S=\PP^2$ mamy $K=-3H$, $H^2=1$,
$\chi(\mathcal O_S)=1$. Dla $D=dH$ otrzymujemy
$\chi(\mathcal O_{\PP^2}(d))=1+\frac12d(d+3)
=\frac{(d+1)(d+2)}2$.
W szczególności dla $d\ge0$ jest to
$\binom{d+2}{2}$, liczba jednorodnych jednomianów
stopnia $d$ w trzech zmiennych.
Dla $d<0$ nie można utożsamiać tej charakterystyki
z wymiarem $H^0$: na przykład
$\chi(\mathcal O(-3))=1$, choć $H^0(\mathcal O(-3))=0$.
Brakująca jedynka pochodzi z $H^2(\mathcal O(-3))$,
co wynika z dualności Serre'a.
\end{przyklad}

\begin{przyklad}[Riemann--Roch na $\PP^1\times\PP^1$]
Niech $S=\PP^1\times\PP^1$ i
$D=aF_1+bF_2$. Wówczas $D^2=2ab$, a
$K=-2F_1-2F_2$, zatem
$D\cdot(D-K)=2ab+2a+2b$.
Ponieważ $\chi(\mathcal O_S)=1$,
otrzymujemy
$\chi(\mathcal O_S(a,b))=(a+1)(b+1)$
dla \emph{wszystkich} liczb całkowitych $a,b$.
Dla $a,b\ge0$ wynika to także bezpośrednio z liczby
bihomogenicznych jednomianów bitypu $(a,b)$.
\end{przyklad}

\section{Hirzebruch--Riemann--Roch i wzór Noether}
\label{sec:ga7-noether}

Rozszerzenie Riemanna--Rocha z samych dzielników
na dowolne wiązki wektorowe wymaga \emph{charakteru Cherna}
i \emph{klasy Todda}. Te obiekty wygodnie definiować
w pierścieniu $\operatorname{CH}^*(X)\otimes_\Z\Q$:
niektóre współczynniki są ułamkowe, lecz po wzięciu
stopnia otrzymujemy liczbę całkowitą.

\begin{definicja}[Charakter Cherna i klasa Todda]
\label{def:ga7-ch-td}
Niech $X$ będzie gładką quasi-rzutową rozmaitością,
$\mathcal E$ wiązką wektorową rzędu $r$.
Po zastosowaniu zasady rozszczepienia niech
$x_1,\ldots,x_r$ oznaczają \emph{formalne pierwiastki Cherna},
tj. pierwsze klasy Cherna kolejnych ilorazów liniowych
w odpowiedniej filtracji pullbacku $\mathcal E$.
Definiujemy
\[
\operatorname{ch}(\mathcal E)=\sum_{i=1}^r e^{x_i},
\qquad
\operatorname{td}(\mathcal E)=\prod_{i=1}^r
\frac{x_i}{1-e^{-x_i}},
\]
przy czym szeregi rozwijamy formalnie i obcinamy
powyżej wymiaru $X$. Dla $x_i=0$ drugi czynnik
rozumiemy jako szereg o wyrazie stałym $1$.
Ponieważ są to wyrażenia symetryczne w $x_i$,
można je zapisać przez klasy Cherna, a więc
zależność od wyboru rozszczepienia znika.
\end{definicja}

\begin{lemat}[Rozwinięcia do stopnia dwa]
\label{lem:ga7-ch-td-surface}
Niech $X$ będzie gładką powierzchnią,
$\mathcal E$ wiązką wektorową rzędu $r$,
a $T_X$ jej wiązką styczną. Wtedy
\[
\begin{aligned}
\operatorname{ch}(\mathcal E)&=r+c_1(\mathcal E)
+\tfrac12(c_1(\mathcal E)^2-2c_2(\mathcal E)),\\
\operatorname{td}(T_X)&=1+\tfrac12 c_1(T_X)
+\tfrac1{12}(c_1(T_X)^2+c_2(T_X)).
\end{aligned}
\]
\end{lemat}
\begin{proof}
Z rozwinięcia $e^x=1+x+x^2/2+\cdots$ otrzymujemy
$\sum e^{x_i}=r+\sum x_i+\frac12\sum x_i^2$.
Podstawiamy
$c_1=\sum x_i$, $c_2=\sum_{i<j}x_ix_j$ i
$\sum x_i^2=c_1^2-2c_2$.
Ponadto zwykły rachunek formalnych szeregów daje
$x/(1-e^{-x})=1+x/2+x^2/12+\cdots$.
Na powierzchni $T_X$ ma rząd dwa, więc mnożąc
czynniki dla jego dwóch pierwiastków
$x_1,x_2$, otrzymujemy wyraz stopnia dwa
$(x_1^2+x_2^2)/12+x_1x_2/4
=(c_1^2-2c_2)/12+c_2/4=(c_1^2+c_2)/12$.
To dowodzi obu wzorów.
\end{proof}

\begin{twierdzenie}[Hirzebruch--Riemann--Roch — twierdzenie zewnętrzne]
\label{tw:ga7-hrr}
Niech $X$ będzie gładką właściwą rozmaitością nad
ciałem $k$ charakterystyki zero, a $\mathcal E$
wiązką wektorową na $X$. Wtedy
\[
\chi(X,\mathcal E)=\deg\bigl(
\operatorname{ch}(\mathcal E)\operatorname{td}(T_X)
\cap[X]\bigr)_{0},
\]
gdzie indeks $0$ oznacza składową cyklu wymiaru zero.
\end{twierdzenie}

\begin{uwaga}[Dlaczego nie przedstawiamy tego jako prostego wniosku?]
Pełny dowód tego twierdzenia wykorzystuje
teorię $K$-grup, deformację do stożka normalnego
albo teorię indeksową. Jest istotnie
głębszy niż dowód wzoru Riemanna--Rocha dla dzielników
na powierzchni, który podaliśmy wyżej.
Przyjmujemy go jawnie jako \emph{wynik zewnętrzny},
a wszystkie poniższe wyprowadzenia są już rachunkowe.
\end{uwaga}

\begin{twierdzenie}[Wzór Noether i uogólniony Riemann--Roch dla powierzchni]
\label{tw:ga7-noether}
Niech $S$ będzie gładką spójną rzutową powierzchnią
nad ciałem algebraicznie domkniętym charakterystyki zero.
Wówczas
\[
12\chi(\mathcal O_S)=K_S^2+\deg c_2(T_S).
\]
Dla wiązki wektorowej $\mathcal E$ rzędu $r$ na $S$
zachodzi ponadto
\[
\chi(\mathcal E)=r\chi(\mathcal O_S)
+\frac12\deg\bigl(c_1(\mathcal E)
(c_1(\mathcal E)-c_1(\omega_S))\bigr)
-\deg c_2(\mathcal E).
\]
\end{twierdzenie}
\begin{proof}
Zastosujmy twierdzenie~\ref{tw:ga7-hrr} do
$\mathcal E=\mathcal O_S$. Wtedy
$\operatorname{ch}(\mathcal O_S)=1$,
a składową stopnia dwa klasy Todda jest
$(c_1(T_S)^2+c_2(T_S))/12$.
Ponieważ $c_1(T_S)=-c_1(\omega_S)$, jej kwadrat
ma stopień $K_S^2$. Otrzymujemy wzór Noether.

Dla ogólnej $\mathcal E$ mnożymy rozwinięcia
z lematu~\ref{lem:ga7-ch-td-surface}. Składowa
stopnia dwa jest równa
\[
\frac r{12}(c_1(T_S)^2+c_2(T_S))
+\frac12c_1(\mathcal E)c_1(T_S)
+\frac12(c_1(\mathcal E)^2-2c_2(\mathcal E)).
\]
Po wzięciu stopnia pierwszy składnik jest równy
$r\chi(\mathcal O_S)$ na mocy dopiero co
wyprowadzonego wzoru, a
$c_1(T_S)=-c_1(\omega_S)$. Pozostałe wyrazy
są iloczynami klas w $\operatorname{CH}^*(S)$ i dają
drugą formułę, ponieważ $c_1(\omega_S)=[K_S]$.
Dla dzielników używamy więc liczby przecięcia,
a dla wiązek ich pierwszych klas Cherna.
Dla wiązki liniowej $\mathcal O_S(D)$ ma ona
$r=1$, $c_1=D$ i $c_2=0$, czyli dokładnie
wcześniej udowodniony wzór powierzchniowy.
\end{proof}

\begin{przyklad}[Charakterystyka holomorficzna powierzchni stopnia $d$]
Niech $S\subseteq\PP^3$ będzie gładką powierzchnią
stopnia $d$. Ze wzorów dla hiperpowierzchni oraz
wzoru Noether otrzymujemy
\[
\chi(\mathcal O_S)=\frac{d(d-4)^2+d(d^2-4d+6)}{12}
=\frac{d(d^2-6d+11)}6.
\]
Dla $d=1,2,3,4$ daje to odpowiednio $1,1,1,2$.
Z drugiej strony dla kwartyki wiemy z rozdziału 6,
że $q=0$, $p_g=1$, więc $\chi=1-q+p_g=2$.
Otrzymaliśmy niezależne sprawdzenie liczby $24$
we wzorze $12\chi=K^2+c_2$.
\end{przyklad}

\begin{przyklad}[Jak zmienia się $c_2$ po rozdmuchaniu?]
\label{prz:ga7-blowup-c2}
Niech $k$ będzie ciałem algebraicznie domkniętym
charakterystyki zero, $S$ gładką rzutową powierzchnią
nad $k$, a $p\in S$ punktem zamkniętym.
Niech $\pi:\widetilde S=\operatorname{Bl}_pS\to S$
będzie jej rozdmuchaniem w $p$.
W rozdziale 6 udowodniliśmy
$K_{\widetilde S}^2=K_S^2-1$.
Standardowy wynik o kohomologii rozdmuchania mówi, że
$\pi_*\mathcal O_{\widetilde S}=\mathcal O_S$ oraz
$R^i\pi_*\mathcal O_{\widetilde S}=0$ dla $i>0$.
Z ciągu spektralnego Leraya wynika wtedy
$\chi(\mathcal O_{\widetilde S})=\chi(\mathcal O_S)$.
Wzór Noether dla obu powierzchni daje
$\deg c_2(T_{\widetilde S})=\deg c_2(T_S)+1$.
Dla rozdmuchania $\PP^2$ w sześciu punktach otrzymujemy
$K^2=9-6=3$ i $\deg c_2=3+6=9$, zgodnie
z niezależnym rachunkiem dla gładkiej powierzchni kubicznej.
W tym argumencie znikanie wyższych obrazów prostych
jest jawną przesłanką zewnętrzną, a nie pominiętym krokiem.
\end{przyklad}

\section{Podsumowanie: dwa poziomy teorii przecięć}
\label{sec:ga7-summary}

W przypadku powierzchni wiele wzorów liczbowych można
wyprowadzić z normalnej wiązki krzywej, adjunkcji
i Riemanna--Rocha na krzywych. To jest \emph{poziom
konkretny}: pozwolił udowodnić Riemanna--Rocha dla
wszystkich dzielników na gładkiej powierzchni.

W wyższych wymiarach konieczne są pełne grupy Chowa,
klasy Cherna, wzory projekcji, przecięcia Gysina
i twierdzenia o przesuwaniu cykli. To jest
\emph{poziom funktorialny}. Zbudowaliśmy jego definicje,
przeprowadziliśmy liczne rachunki i wskazaliśmy,
które twierdzenia konstrukcyjne nadal są przyjęte
z zewnątrz. W ten sposób teoria z rozdziału 6
otrzymała precyzyjne podstawy, a wzór Noether
stał się konkretną konsekwencją
Hirzebrucha--Riemanna--Rocha.

\begin{przyklad}[Wielomian Hilberta powierzchni jako kontrola wzoru Riemanna--Rocha]
\label{prz:ga7-hilbert-polynomial}
Niech $S\subseteq\PP^3_k$ będzie gładką powierzchnią
stopnia $d\ge1$ nad ciałem algebraicznie domkniętym
charakterystyki zero, a $H$ dzielnikiem hiperpłaszczyznowym
na $S$. Oznaczmy $h=[H]=c_1(\mathcal O_S(1))$.
Dla dowolnego $m\in\Z$ wzór Riemanna--Rocha
i równości $K_S\sim(d-4)H$, $\deg(h^2)=d$,
gdzie $H$ jest dzielnikiem hiperpłaszczyznowym
oraz $h=[H]\in\operatorname{CH}^1(S)$, dają
\[
\chi(\mathcal O_S(m))
=\frac{d(d^2-6d+11)}6
 +\frac d2\,m(m+4-d).
\]
Rachunek ten można niezależnie sprawdzić za pomocą
dokładnego ciągu snopów na $\PP^3$:
\[
0\longrightarrow\mathcal O_{\PP^3}(m-d)
\longrightarrow\mathcal O_{\PP^3}(m)
\longrightarrow\mathcal O_S(m)\longrightarrow0.
\]
Ponieważ
$\chi(\mathcal O_{\PP^3}(t))=\binom{t+3}{3}$
dla każdego $t\in\Z$, z addytywności charakterystyki
Eulera otrzymujemy również
\[
\chi(\mathcal O_S(m))
=\binom{m+3}{3}-\binom{m-d+3}{3}.
\]
W tym wzorze symbol Newtona oznacza wielomian
$\binom{t}{3}=t(t-1)(t-2)/6$, także dla ujemnego
$t$. Rozwinięcie obu wyrażeń daje tę samą funkcję
kwadratową zmiennej $m$. Dla $m=1$ gładka
powierzchnia kubiczna ($d=3$) oraz gładka kwartyka
($d=4$) mają obie $\chi(\mathcal O_S(1))=4$,
choć ich charakterystyki strukturalne wynoszą
odpowiednio $1$ i $2$. Pojedyncza wartość
wielomianu Hilberta nie rozstrzyga więc o typie
biwymiernym powierzchni.
\end{przyklad}

\begin{uwaga}[Przygotowanie do geometrii biwymiernej]
W rozdziale 8 liczby $C^2$, $K_S\cdot C$ i krzywe
z ujemnym samoprzecięciem będą kontrolować możliwe
kontrakcje. Program modeli minimalnych w wymiarze dwa
korzysta szczególnie z tego, że przecięcia z dzielnikiem
kanonicznym rozpoznają krzywe wyjątkowe. Teoria
klas Cherna dostarczy z kolei niezmienników powierzchni,
które pozwalają odróżniać typy biwymierne.
\end{uwaga}

% Źródła matematyczne użyte do weryfikacji:
% Stacks Project: Chow Homology and Chern Classes (tagi 02U4, 02TZ, 02UF),
% Grothendieck-Riemann-Roch (odpowiedni rozdział Stacks Project).
% Fulton, Intersection Theory; Hartshorne, Algebraic Geometry V.
% Wykorzystane zewnętrznie: moving lemma, lokalizacja i homotopia CH,
% projective bundle formula, Whitney/naturalność, sekcje regularne,
% HRR, Bertini i obszerność, wyższe obrazy proste rozdmuchania.

% =============================================================
% ROZDZIAL 8: GEOMETRIA BIWYMIERNA I MODELE MINIMALNE
% =============================================================
